\section{Основные определения}

Введём основные определения, необходимые для формализации задачи:

\begin{itemize}
    \item Граф \(G = (V, E)\) — ориентированный граф, где \(V\) — множество вершин (\(|V| = n\)), а \(E\) — множество рёбер.
    \item Вес ребра \(w(u, v)\) для ребра \((u, v) \in E\) определяется функцией \(w : E \to \mathbb{R}_{\geq 0}\), где \(w(u, v) \geq 0\) для всех рёбер.
    \item Кратчайший путь из вершины \(s \in V\) в вершину \(v \in V\) — это путь \(P \subseteq E\), минимизирующий сумму весов рёбер:
    \[
    W(P) = \sum_{(u, v) \in P} w(u, v).
    \]
    \item Минимальное расстояние от вершины \(s\) до вершины \(v\) обозначается \(d(s, v)\) и определяется как:
    \[
    d(s, v) = \min_{P \in \mathcal{P}(s, v)} W(P),
    \]
    где \(\mathcal{P}(s, v)\) — множество всех путей от \(s\) до \(v\).
    \item Если вершина \(v\) недостижима из вершины \(s\), то \(d(s, v) = \infty\).
\end{itemize}